\section{人工智能更像制造业}

Peak 说 AI 更像制造业。这句话很反直觉，但非常有解释力。很多人把 AI 想成软件：模型训练好，产品上线，边际成本很低。但 Agent 公司面对的实际问题更像制造业：质量、成本、产能、供应链、交付、维护、可复现和客户场景适配。

\begin{figure}[H]
\centering
\includegraphics[width=0.96\textwidth]{ai-like-manufacturing.png}
\caption{AI 更像制造业：模型、算力、数据、交付和质量。}
\end{figure}

\begin{knowledgebox}{读图：制造业隐喻讲的是可交付性}
图中模型只是底座，算力是产能，数据是原料和工艺，交付是客户现场，质量是稳定可复现。Agent 公司要把这些都压实，而不是只展示模型能力。
\end{knowledgebox}

\subsection{质量、成本和产能}

制造业关心良率、成本、产能和供应链。AI Agent 也类似：同一个任务能否稳定完成，失败率能否降低，推理成本是否可控，工具调用是否可靠，权限和支付是否打通，环境阻塞是否减少。这些问题看起来不酷，却决定用户能不能依赖产品。

\begin{warningbox}{不要把 Agent 只当软件功能}
如果一个 Agent 每次成功率不稳定、成本不可控、环境经常阻塞，它就更像一个手工 demo，而不是可交付产品。制造业视角能提醒团队关注稳定性和良率。
\end{warningbox}

\subsection{生态配套：支付、权限和反爬}

Peak 提到 Cloudflare、人机验证、网站不让 Agent 访问，以及 Stripe 做 Agent 支付。这说明 Agent 不只是模型问题，还需要整个互联网生态适配：身份、支付、权限、反爬、审计、责任。没有这些配套，Agent 很多能力会卡在真实世界边界。

\subsection{Agent 基础设施的四个缺口}

前面提到支付、权限和反爬只是例子，本节把这些缺口系统化。如果把 Agent 当成一个可以替你上网、下单、查资料、操作工具的数字员工，它需要的基础设施就不只是模型 API。首先是身份和权限：Agent 代表谁行动，能访问哪些网站，能不能保存登录态。其次是支付：Agent 如果要买东西、订服务、付费调用工具，就需要机器可用的支付接口。第三是反爬和人机验证：很多网站默认把自动化访问视为风险。第四是审计和责任：Agent 做错了，谁负责，如何回滚，如何追踪操作链。

\begin{knowledgebox}{Agent 基础设施清单}
\begin{tabularx}{\textwidth}{p{0.22\textwidth}p{0.34\textwidth}X}
\toprule
基础设施 & 解决的问题 & 如果缺失会怎样 \\
\midrule
身份/权限 & Agent 代表谁、能做什么 & 无法安全进入真实账户和业务系统。 \\
支付 & Agent 能否完成交易和购买 & 只能浏览和建议，难以闭环执行。 \\
反爬协作 & 网站是否允许 Agent 访问 & 经常被 Cloudflare、人机验证或风控拦住。 \\
审计/回滚 & 操作可追踪、错误可纠正 & 企业和高价值任务不敢交给 Agent。 \\
\bottomrule
\end{tabularx}
\end{knowledgebox}

\begin{importantbox}{Agent 爆发不只靠模型}
当模型足够强以后，决定 Agent 能否爆发的瓶颈会转向身份、支付、权限、反爬、审计和责任这些“互联网基础设施”。
\end{importantbox}

\subsection{本章小结}

AI 像制造业，指的是 Agent 产品最终要面对质量、成本、产能和交付问题。模型能力只是开端，稳定进入真实工作流才是关键。

